\section{Related Work}
\label{sec:related}
\paragraph{Privacy and contextual norms in LM agents.}
PrivacyLens evaluate whether language models recognize contextually appropriate disclosure norms~\citep{shao2024privacylens}; AgentDAM extends this lens to autonomous web agents with data minimization~\citep{zharmagambetov2026agentdam}. These treat the agent as a single decision-making unit. We shift the unit of analysis to the \emph{handoff artifact} produced when one agent compresses upstream interaction for another, where a summarizer can preserve the sensitive fact while silently dropping the marker that constrained how it could be used.
 
\paragraph{Privacy leakage in multi-agent systems.}
AgentLeak shows violations frequently occur in internal inter-agent channels rather than user-facing outputs~\citep{el2026agentleak}; SumLeak formalizes \emph{compositional} leakage in which individually innocuous responses aggregate into a sensitive attribute~\citep{patil2025sum}; and prompt injection or environmental manipulation can elicit sensitive context at execution time~\citep{alizadeh2025simple, liao2025eia}. Our setting is non-adversarial: the sensitive information is already present at the source and is often accompanied by an explicit boundary marker. The failure we characterize is that the marker, not the fact, is disproportionately weakened or dropped in the cross-phase compression step. 

\paragraph{Memory and inter-agent communication as channels.}
MEXTRA demonstrates black-box extraction of stored user--agent interactions~\citep{wang2025unveiling}; Prompt Infection~\citep{lee2025prompt}, MASLeak~\citep{wang2025ip}, and Agent-in-the-Middle~\citep{he2025red} show inter-agent channels can be exploited to self-replicate instructions, extract proprietary system information, or compromise an entire MAS. Similarly, \citet{wang2026state} show that persistent agent state can carry influence that a deployed monitor scores as safe, quantified as a sub-threshold propagation gap. We adopt the same channel-level lens but study a non-adversarial failure: when one agent's output becomes another's input via summarization, what governance information survives, what is silently lost, and when does that loss become final audience-facing leakage? 
 
\paragraph{Communication Privacy Management.}
CPM theory models privacy as boundary work: individuals own private information, license its disclosure to specific co-owners under specific use conditions, and experience boundary turbulence when those conditions are violated~\citep{petronio1991communication, petronio2002boundaries, griffin2006first}. 
We follow \citet{petronio2002boundaries} specifically. whose rule-management heuristics (culture, context, motivation, risk--benefit ratio) govern \emph{when} a boundary rule forms; our four marker categories instead capture the \emph{content} such a rule must carry once formed, operationalizing CPM's boundary-rule dimensions of ownership, co-ownership, and permeability for the multi-agent setting.
We instantiate this content operationally rather than formalizing CPM as a whole, which we scope as beyond this paper. 
In CPM terms our contribution is to identify a structural failure of MAS where operational content consistently survives while the licensing rule does not. We name this failure mode \emph{summary collapse}, measure it with $\sigma_b$, $\sigma_{\mathrm{op}}$, and $\rho$, and show through controlled downstream tests that boundary representation determines whether context dump becomes final leakage.
